% ============================================================================
% APPENDIX B — GLOSSARY OF DEEP LEARNING TERMS
% ============================================================================
\chapter{Glossary of Deep Learning Terms}
\label{app:dl_glossary}

This appendix provides concise definitions of deep learning concepts referenced throughout the document, intended as a quick reference for readers with a communications background.

\begin{description}[style=nextline, font=\normalfont\bfseries, leftmargin=1.5em]

      \item[Neuron]
            The basic computational unit of a neural network. It computes a weighted sum of its inputs, adds a bias, and applies an activation function: $y = f(\bm{w}^T \bm{x} + b)$.

      \item[Weight]
            A learnable parameter in a neural network that scales the connection between two neurons. The collection of all weights (and biases) constitutes the model's trainable parameters, denoted $\bm{\theta}$.

      \item[Bias]
            A learnable constant added to the weighted sum in a neuron before the activation function is applied: $z = \bm{W}\bm{x} + \bm{b}$. It allows the network to shift the activation function, giving it more flexibility.

      \item[Activation Function]
            A nonlinear function applied element-wise to the output of a neural network layer. Common choices include ReLU, sigmoid, tanh, softmax, and the linear (identity) function. Without nonlinear activations, a multi-layer network would collapse to a single linear transformation.

      \item[Fully Connected (FC) Layer]
            A layer where every neuron is connected to every neuron in the previous layer via a weight matrix $\bm{W}$. Also called a dense layer. The computational cost of an FC layer scales as $\mathcal{O}(n_{\mathrm{in}} \times n_{\mathrm{out}})$.

      \item[Forward Pass]
            The computation that propagates input data through the network, layer by layer, to produce an output. During inference (deployment), only the forward pass is needed---backpropagation is not required.

      \item[Inference]
            The process of using a trained model to make predictions on new, unseen data. In the PAPR context, inference means performing one forward pass per OFDM symbol at real-time rates.

      \item[Loss Function]
            A scalar function that quantifies how far the network's output is from the desired output. Training minimises this function. Common choices include MSE for regression and cross-entropy for classification. In PAPR, custom losses combine PAPR and MSE terms.

      \item[Training]
            The iterative process of adjusting network parameters to minimize the loss function. Involves repeated forward passes, loss computation, backpropagation, and parameter updates over many epochs.

      \item[Gradient]
            The vector of partial derivatives of the loss function with respect to each trainable parameter. It indicates the direction and magnitude of the steepest increase in loss; the optimizer updates parameters in the opposite direction.

      \item[Backpropagation]
            The algorithm used to compute gradients of the loss function with respect to all network parameters by applying the chain rule of calculus, layer by layer, from the output back to the input. These gradients are then used by the optimizer to update the weights.

      \item[Optimizer]
            The algorithm that updates network parameters based on computed gradients. SGD (Stochastic Gradient Descent) is the simplest; Adam (Adaptive Moment Estimation) is the most commonly used in practice, combining momentum and per-parameter learning rate adaptation.

      \item[Learning Rate ($\eta$)]
            A hyperparameter that controls the step size of each parameter update. Too large a value causes training to diverge; too small a value makes convergence extremely slow.

      \item[Epoch]
            One complete pass through the entire training dataset. Training typically requires hundreds of epochs before the network converges to a good solution.

      \item[Batch]
            A subset of the training data processed together in one forward and backward pass. Using batches (rather than the entire dataset) enables efficient GPU utilization and introduces beneficial stochastic noise into training.

      \item[Convergence]
            The point during training at which the loss function stabilizes and further epochs yield negligible improvement. A model that fails to converge may be underfitting (too simple) or have a learning rate that is too high.

      \item[Hyperparameter]
            A configuration value set before training begins and not learned from data. Examples include learning rate, batch size, number of layers, number of epochs, and the loss balance factor $\alpha$.

      \item[Overfitting]
            When a network learns to memorize the training data rather than generalize to unseen data. Symptoms include very low training loss but high validation/test loss. Mitigated by dropout, regularization, and early stopping.

      \item[Regularization]
            Any technique that prevents overfitting by constraining the model's complexity. Examples include dropout, weight decay ($L_2$ regularization), and batch normalization.

      \item[Dropout]
            A regularization technique that randomly sets a fraction of neuron outputs to zero during each training step. This prevents neurons from co-adapting and encourages the network to learn redundant representations, reducing overfitting.

      \item[Batch Normalization (BatchNorm)]
            A technique that normalizes the inputs to each layer by subtracting the batch mean and dividing by the batch standard deviation. This stabilizes training, allows higher learning rates, and acts as a mild regularizer.

      \item[Vanishing Gradients]
            A problem where gradients become extremely small as they propagate through many layers during backpropagation, causing early layers to learn very slowly or not at all. ReLU activations and residual connections help mitigate this.

      \item[Autoencoder]
            A neural network architecture consisting of an encoder that maps input data to a lower-dimensional (or transformed) representation, and a decoder that reconstructs the original data. In the PAPR context, the encoder produces a low-PAPR signal and the decoder recovers the original data symbols.

\end{description}
